The present results support the use of ERP-window features rather than spectral band-power summaries for RSVP attention detection, because the target-related signal in this paradigm is expected to be dominated by transient posterior ERP activity rather than sustained oscillatory changes. The importance analysis concentrated on the 400-500 ms interval and on posterior channels such as PO4, PO7, O1, which is neurophysiologically consistent with a P300-like target response over parieto-occipital scalp regions. The corrected posterior ROI difference wave also showed its annotated extremum within the target-related window rather than at the earliest boundary, which aligns the waveform interpretation with the channel-time importance summary and with the expected posterior P300 timing. At the same time, the moderate classification performance remains scientifically plausible given the strong class imbalance, sparse eight-channel montage, and cross-subject variability characteristic of rapid visual presentation datasets. These properties make the proposed approach valuable as a lightweight, explainable baseline that captures meaningful target-related ERP structure without relying on heavy preprocessing or opaque deep models.
